\subsection{A common cusp colour across inert primes}

\begin{statement}{B.42}{Cross-prime colour rigidity and its compact limit}{Proved}
Let $E:X^3+Y^3+Z^3=0$, let $D=E[3]$ be its nine coordinate
cusps, and use the unit tangent at $O=C_0=[1:-1:0]$ with
parameter $Z/X$.  Let $N=\pi_1((E-D)_{\overline{\mathbb Q}})^{(3)}$
be the full geometric pro-three group, with its actual based action,
and put $T=T_3E$, $K=\mathbb Q(\zeta_3)$ and
$P=G_{K,\{3\}}^{(3)}$.

Use the compact matrices of B.24 and B.28:
\[
 R=\begin{pmatrix}0&-1\\1&-1\end{pmatrix},\qquad
 C=\begin{pmatrix}1&-1\\0&-1\end{pmatrix},\qquad
 T_\tau=I+3R,\quad T_{\bar\tau}=I+3R^2.
                                                               \tag{B.42.1}
\]
Here $\tau,\bar\tau=c\tau c$ are the actual conjugate Frobenii
above seven forming the free basis of B.28.  Let $R_0\in E(\mathbb Q)[3]$
be the point represented by $e_1/3$ in $T\otimes\mathbb Q_3/T$.
Define the compact cocycle $\beta$ by
\[
 \beta(c)=0,\qquad \beta(\tau)=a=(1,2)^t,
 \qquad\beta(\bar\tau)=-a.
                                                               \tag{B.42.2}
\]
For every actual full $N$-section, its compact class has a unique form
\[
 \alpha=b\beta+k\kappa_3(R_0),\qquad
 b\in\mathbb Z_3,\quad k\in\mathbb Z/3.
                                                               \tag{B.42.3}
\]
At every inert rational prime $q\ne3$,
\[
 \operatorname{loc}_q\beta
 \in27H^1_{\mathrm{unr}}(\mathbb Q_q,T),\qquad
 H^1_{\mathrm{unr}}(\mathbb Q_q,T)\simeq
       \mathbb Z/3^{v_3(q+1)}.
                                                               \tag{B.42.4}
\]
Consequently every ramified inert prime satisfying
\[
                 v_3(q+1)\le3+v_3(b)                         \tag{B.42.5}
\]
has the same rational cusp colour $O+kR_0$.  We set $v_3(0)=+\infty$.
In particular the colours agree at every two ramified inert primes
with $q\not\equiv-1\pmod{81}$.  If $b=0$, all ramified good
primes, split or inert, have colour $O+kR_0$.

At a split prime $q\equiv4,7\pmod9$, choose a place of $K$
above $q$, and write its Frobenius class in $P/\Phi(P)$ as
$u\tau+v\bar\tau$.  Put $P_0=(e_1+2e_2)/3\bmod T$.
Then $u+v\ne0$ in $\mathbb F_3$, and a ramified section's
colour necessarily satisfies
\[
 \text{centre}-O=kR_0+(b\bmod3)\frac{u-v}{u+v}P_0.
\]
For $\alpha=\beta$, this label is $P_0$ at seven with
$\zeta_3\equiv2$, and zero at sixty-one.  This is a necessary
compact constraint; a full section ramified at both primes has
not been constructed.

The factor twenty-seven is sharp for the compact necessary condition:
for each $r\in\{0,1,2\}$ a positive-density set of inert primes
$q\equiv161\pmod{243}$ satisfies
$\operatorname{loc}_q\beta=r\kappa_q(R_0)$.
The class $\beta$ itself is the compact image of actual full
$N$-sections unramified outside three.  These compact labels do not
assert local cuspidality or create ramification at those primes.
Thus the unrestricted assertion B.10(a) remains open.
\end{statement}
\begin{proof}
\emph{The whole compact cocycle and its two coordinates.}
At $q\ne3$ the coefficient action on $T$ is unramified.
Wild inertia is pro-$q$ and maps trivially to $T$; the remaining
inertia map factors through $\mathbb Z_3(1)$.  Its value $v$
satisfies $F_qv=qv$.  Compact Frobenius eigenvalues have absolute
value $\sqrt q$, so $F_q-q$ is rationally invertible and $v=0$.
Thus every compact cocycle, including one obtained from an infinitely
ramified full section, is wholly unramified outside three.

Over $K=\mathbb Q(\zeta_3)$ the coefficient action is pro-three.
The action/cocycle graph factors through
$P=G_{K,\{3\}}^{(3)}$, and the relative arithmetic group is
$\Gamma=P\rtimes\langle c\rangle$.  B.28 proves that
$\tau,\bar\tau$ are a free basis of $P$.  Two is invertible, so
every compact cocycle can be normalized to vanish at $c$.
It is then determined by its value $w=(w_1,w_2)^t$ at $\tau$;
its value at $\bar\tau$ is $Cw$.
The remaining coboundaries have conjugator in
$T^c=\mathbb Z_3e_1$ and change $w$ by a multiple of
$(T_\tau-I)e_1=3e_2$.  The rational torsion point $R_0$ has
the actual Kummer cocycle
\[
 \kappa_3(R_0)(g)=(g-1)e_1/3,
 \qquad \kappa_3(R_0)(\tau)=e_2,
 \qquad \kappa_3(R_0)(c)=0.
\]
This proves (B.42.3), with
\[
                   b=w_1,\qquad k=w_2-2w_1\pmod3.
                                                               \tag{B.42.6}
\]
No passage from an arbitrary compact class to a ramified full
section is used here.

\emph{The exact compact affine image.}
Use the product $(z,M)(w,L)=(z+Mw,ML)$.
The graph image $H$ of $\beta$ is generated topologically by
\[
 A=(a,T_\tau),\qquad B=(-a,T_{\bar\tau}),\qquad c=(0,C).
\]
Its pro-three subgroup $H_K$ is generated by $A,B$.
Since $T_\tau T_{\bar\tau}=7I$ and
$T_\tau+T_{\bar\tau}=-I$, one obtains
\[
 [A,B]=(3a,I),\qquad
 L=3\mathbb Z_3a+9\mathbb Z_3Ra
       =3\mathbb Z_3a\oplus27\mathbb Z_3e_1.
                                                               \tag{B.42.7}
\]
The normal closure of this commutator contains $9Ra$, by
conjugating $3a$ with $A$ and subtracting $3a$.
The displayed lattice is stable under $T_\tau,T_{\bar\tau}$,
their inverses, and $C$.  It is therefore the closed normal
translation subgroup generated by the commutator.

It is also the entire coefficient kernel.  Indeed the logarithms
of $1+3\zeta_3$ and $1+3\zeta_3^2$, divided by three,
reduce to the basis $\zeta_3,\zeta_3^2$ modulo three.
Consequently their powers identify $\mathbb Z_3^2$ with
$1+3\mathcal O_{K,3}$.  The quotient $H_K/L$ is abelian,
generated by the images of $A,B$, and the composite
\[
 \mathbb Z_3^2\longrightarrow H_K/L
       \longrightarrow1+3\mathcal O_{K,3}
\]
is an isomorphism.  The first arrow is surjective; therefore
both arrows are isomorphisms.  This proves the claimed entire
kernel, including its closed inverse-limit topology.  In particular
\[
                         L^c=27\mathbb Z_3e_1.              \tag{B.42.8}
\]

The translation $-3a/2$ lies in $L$, and direct multiplication gives
\[
 D_0=(-3a/2,I)AB=(9e_1/2,7I)\in H_K^c.
\]
Seven topologically generates $1+3\mathbb Z_3$.
For $t=7^u$, continuous powering gives the actual lift
\[
                  D_0^u=(3(t-1)e_1/4,tI).                 \tag{B.42.9}
\]
Every $c$-fixed element with coefficient $tI$ differs from
(B.42.9) by (B.42.8).  Its translation is therefore
$3(t-1)e_1/4+27ze_1$ for some $z\in\mathbb Z_3$.

\emph{Inert Frobenius, with its involution and integral conjugator.}
Let $g\in H$ be the actual affine Frobenius at an inert prime $q$.
Its closed procyclic group has a unique order-two component $d$.
Write $d=hc$, with $h\in H_K$.  Then $c(h)=h^{-1}$, and the
unique square root $u=h^{1/2}$ in its pro-three procyclic subgroup
satisfies $ucu^{-1}=d$.  Conjugating by $u$ makes the order-two
component literally $c$.  The resulting Frobenius is $h'c$ with
$h'\in H_K^c$.

The centralizer of $C$ in the CM coefficient algebra consists of
scalars.  Thus the coefficient of $h'$ is the unique
$t\in1+3\mathbb Z_3$ with $t^2=-q$: inert elliptic Frobenius
has square $-qI$.  Its translation has the form just obtained.
On $T^c$ the local coboundary operator is $t-1$, whereas on the
negative eigenspace it is $-t-1$, a unit.  The first translation
term in (B.42.9) is the integral local coboundary
$(t-1)(3e_1/4)$, leaving $27ze_1$.
Moreover
\[
                       v_3(t-1)=v_3(q+1),
\]
since $t+1$ is a unit.  This proves (B.42.4).

An affine conjugator $(v,J)\in H_K$ replaces $(z,M)$ by
\[
 \bigl(J^{-1}(z+(M-1)v),\ J^{-1}MJ\bigr).
\]
It preserves the divisibility assertion in integral local
cohomology.  Since $J\equiv I\pmod3$, it also fixes
$e_1/3\bmod T$, and hence preserves the comparison with the
specified rational Kummer class.  This checks that the preceding
normalization has not changed the colour coordinate.

Combining (B.42.3)--(B.42.4) gives
$\operatorname{loc}_q\alpha=\kappa_q(kR_0)$ under (B.42.5).
At inert $q$ the only rational local cusps are the three rational
cusps.  Their differences are distinct three-torsion points and
remain distinct under good reduction and local completed Kummer
theory.  Thus a locally cuspidal section has exactly the stated
centre.  If $b=0$, the same injectivity on all nine local
three-torsion points applies at split primes as well.

\emph{The precise ramified-to-cuspidal implication used here.}
The finite-cover argument of A.28 applies directly to the actual
geometrically pro-three extension at $q\ne3$; it does not require
that an $N$-section lift to a full profinite section.  Here are its
data in this case.  The good model and disjoint cusp divisor give
the based unramified specialization for finite three-groups
\citep[XIII, Corollary 2.9 and construction 2.10]{BSGA1}.
Each cusp inertia is a cyclic free pro-three factor, and distinct
cusps can occur together in a free basis.  The centralizer and
malnormality results of \citet[Theorems A' and B']{cHerfortRibes1985}
therefore apply.

If a nonzero tame geometric inertia group $I\simeq\mathbb Z_3$
lay in no cusp factor, its primitive annulus $I\setminus I^3$
would be disjoint from the compact union of all conjugate cusp
groups.  Indeed any nontrivial intersection would put $I$ in
that factor, by the centralizer and root property of a cyclic
free factor.  Separate the compact sets in a finite quotient,
with characteristic open kernel $M$, refined so that
$I\cap M\subset I^3$; this also follows from the separation
because the identity belongs to every cusp group.
The open group $MI$, stable under the section's
Frobenius $\psi$, has a quotient character
$MI\to I/I^3$ detecting $I$ and killing every cusp inertia of
this cover.  The arithmetic group generated by $MI$ and $\psi$
defines a geometrically connected finite cover of the actual
special fibre; the map from $\widehat{\mathbb Z}$ sending one to
$\psi$ gives its descent.  Pro-three covers of this cover have
pro-three Galois closure over the original curve by the finite
wreath-product construction.  Thus the geometric pro-three group
really is $MI$.

Fill the punctures and abelianize.  The detecting character factors
through the torsion-free module $T_3\operatorname{Jac}(Y)$ of the
compactification.  Its nonzero inertia image is multiplied by $q$
under Frobenius, contradicting the eigenvalue absolute value
$\sqrt q$ \citep[II, Theorem 1.1(b); III, Proposition 9.1]{Milne2008}.
Hence $I$ lies in a unique cusp factor.  The full local section
normalizes this factor and lies in its decomposition group;
surjectivity to local Galois makes the cusp $\mathbb Q_q$-rational.
This recovers the good-place full-$N$ implication of LM7 E.61
with its required hypotheses verified.  It makes no inference
about arbitrary quotients from their abelian inertia alone.

\emph{An elementary check of the colour coordinate.}
For $q\equiv2,5\pmod9$, put $e_q=(q^2-1)/3\pmod3$.
The inert Frobenius row on $(\zeta_3,1-\zeta_3)$ is
$(e_q,2e_q)$: use $(1-\zeta_3)^2=-3\zeta_3$ and the fact
that every element of $\mathbb F_q^\times$ is a cube.
Against the seven-adic rows $(2,0),(2,2)$ this gives
$F_q^2=e_q(\tau+\bar\tau)$ in $P/\Phi(P)$.
Since the compact action is trivial modulo three over $K$,
\[
 \overline\alpha(F_q^2)=e_q(w+Cw)
                     =-e_q(w_1+w_2)e_1.
\]
The corresponding value of $\kappa(R_0)$ is $-e_qe_1$.
The local compact group has order three, and restriction to the
unramified quadratic extension is injective because corestriction
composed with restriction is multiplication by two.  Thus the
direct residue calculation gives $k=w_1+w_2\pmod3$, agreeing
with (B.42.6).  The independent program
\path{c_inert_cusp_rigidity.py} checks all compatible mod-nine
compact matrices and the literal cubic residue rows.

\emph{Sharpness of the compact bound, without asserted ramification.}
Set $t=7^{27}$, so $v_3(t-1)=4$.  For each $r\in\{0,1,2\}$,
the exact image $H$ contains
\[
       g_r=\bigl(r(t-1)e_1/3,I\bigr)D_0^{27}c,
\]
because $(t-1)/3\in27\mathbb Z_3$.  Its coefficient is $F=tC$,
and its translation differs from $r(F-1)e_1/3$ by the integral
coboundary $(F-1)(3e_1/4)$.  Modulo $243$ the data are
\[
 F=\begin{pmatrix}163&80\\0&80\end{pmatrix},\qquad
 z_r=\binom{54r}{0},\qquad \det F=161\pmod{243}.
                                                               \tag{B.42.10}
\]
Apply Chebotarev to the actual finite action/cocycle quotient
modulo $243$ and the conjugacy class of this element
\citep[V, Theorem 3.23; VIII, Theorem 7.4]{BMilneCFT2020}.
It gives a positive-density set of inert rational primes in the
stated congruence class.  The congruence determines the integral
compact class: $F_q-1$ has Smith valuations zero and four, so
$3^4T\subset(F_q-1)T$.  Perturbing $F$ modulo $3^5$ changes
$(F-1)e_1/3$ only by $3^4T$; the cocycle perturbation is in
$3^5T$.  The asserted equality modulo local coboundaries follows.
It is invariant on the whole conjugacy class because $R_0$ is
rational.

Finally choose any $n\in N$ mapping to $(1,2)$ in $T$.
The actual pro-three action and the free basis of $P$ allow graph
values $n,c(n)$ on $\tau,\bar\tau$, with value one at $c$.
The free universal property gives an actual continuous full
$N$-section unramified outside three whose compact class is
$\beta$.  Thus the three compact Kummer labels genuinely arise
from a compact class of a full section.  The argument supplies
neither nonzero inertia nor full local cusp geometry at the
Chebotarev primes, and hence no two-prime counterexample to B.10(a).

The independent exact program \path{c_inert_affine_colours.py}
enumerates all $118{,}098$ elements of the affine image modulo
$81$, verifies the kernel and its $c$-fixed line, and checks every
inert element with valuation at most three.  It also verifies
(B.42.10).  The infinite-level kernel is proved by (B.42.7)
and the independent CM logarithms, not inferred from the finite
enumeration.

\emph{The split-prime necessary constraint.}
For $q\equiv4,7\pmod9$, write its chosen $K$-Frobenius class
as $(u,v)$ in the actual basis of $P/\Phi(P)$.  Modulo nine,
\[
 F_q-I=3D_q,\qquad D_q=uR+vR^2\pmod3,
 \qquad\det D_q=(u+v)^2\ne0.                              \tag{B.42.11}
\]
The inequality follows either from $\det F_q=q$ and
$\det T_\tau=\det T_{\bar\tau}=7$, or from
$q\not\equiv1\pmod9$.  Hence $(F_q-I)T=3T$, and the
whole local compact group is $E[3]$.  The Kummer class of
$x/3\bmod T$ is represented modulo three by $D_qx$.
Because the action on $T/3T$ is trivial over $K$,
\[
 \alpha(F_q)=b(u-v)a+kD_qe_1\pmod3,
 \qquad D_qa=(u+v)a\pmod3.
\]
Inverting $D_q$ proves the stated colour formula.
This calculation uses only the Frobenius class in the Frattini
quotient; it does not replace an actual Frobenius word by a word
in the free generators at any higher level.

For an explicit residue test set
$i=\kappa(\zeta_3)(F_q)$ and
$j=\kappa(1-\zeta_3)(F_q)$ in $\mathbb F_3$.
The actual seven-adic Kummer rows give
\[
 i=2(u+v)\ne0,\qquad j=2v,
 \qquad\frac{u-v}{u+v}=1+\frac ji.                       \tag{B.42.12}
\]
At seven with $\zeta_3\equiv2$, the row is $(2,0)$.
At sixty-one the two roots $13,47$ both give the row $(2,1)$.
Thus the respective coefficients of $P_0$ are one and zero.
Complex conjugation negates both $P_0$ and this coefficient,
so the local point is independent of how the conjugate place is
labelled.  The independent program
\path{c_split_compact_colours.py} checks the matrix identity
for every relevant mod-three cocycle and the literal residue
rows for all such split primes below one thousand.
Neither these compact constraints nor the preceding sharpness
construction settles the nonlinear compatibility required for
nonzero inertia at two rational primes.
\end{proof}
